
% ------------------------------------------------------------
\begin{frame}{Intervalos de confianza}

Un intervalo de confianza proporciona un rango de valores
compatibles con los datos y el procedimiento estadístico.

Ejemplo:

\[
\widehat{\Delta}=2.4\text{ cm}
\]

\[
IC_{95\%}=[0.8,\;4.0]\text{ cm}
\]

Podemos comunicar:

\begin{quote}
La diferencia estimada fue de 2.4 cm, con un intervalo de confianza
del 95\% de 0.8 a 4.0 cm.
\end{quote}

\end{frame}

% ------------------------------------------------------------
\begin{frame}{Pruebas de hipótesis comunes}

\begin{center}
\begin{tabular}{lll}
\toprule
\textbf{Pregunta} & \textbf{Prueba frecuente} & \textbf{Alternativa}\\
\midrule
Una media vs. valor & $t$ de una muestra & Wilcoxon\\
Dos grupos & $t$ de dos muestras & Mann--Whitney\\
Antes/después & $t$ pareada & Wilcoxon pareada\\
Más de dos grupos & ANOVA & Kruskal--Wallis\\
Dos variables categóricas & $\chi^2$ & Fisher\\
Asociación lineal & correlación $t$ & Spearman\\
\bottomrule
\end{tabular}
\end{center}

\end{frame}

% ------------------------------------------------------------
\begin{frame}{ANOVA: una advertencia importante}

Un resultado significativo de ANOVA nos dice:

\[
\boxed{\text{no todas las medias son iguales}}
\]

Pero no identifica automáticamente:

\[
\text{qué grupos difieren entre sí}
\]

Para eso se requieren comparaciones posteriores apropiadas
(post-hoc), teniendo en cuenta la multiplicidad.

\end{frame}

% ------------------------------------------------------------
\begin{frame}{Pruebas no paramétricas}

Cuando las condiciones de una prueba paramétrica no son razonables,
pueden existir alternativas.

Ejemplos:

\begin{center}
\begin{tabular}{ll}
\toprule
Paramétrica & Alternativa frecuente\\
\midrule
$t$ independiente & Mann--Whitney\\
$t$ pareada & Wilcoxon\\
ANOVA & Kruskal--Wallis\\
Correlación de Pearson & Spearman\\
\bottomrule
\end{tabular}
\end{center}

\begin{alertblock}{Pero}
``No paramétrica'' no significa ``sin supuestos''.
\end{alertblock}

\end{frame}

% ------------------------------------------------------------
\begin{frame}{Supuestos: ¿por qué importan?}

Una prueba estadística depende de un modelo.

Por ejemplo, para una prueba $t$ pueden ser importantes:

\begin{itemize}
    \item independencia de las observaciones,
    \item estructura adecuada del diseño,
    \item distribución aproximadamente apropiada de los residuos,
    \item ausencia de valores extremos que dominen el resultado.
\end{itemize}

\begin{block}{Principio}
La prueba debe corresponder al diseño y a la estructura de los datos.
\end{block}

\end{frame}

% ------------------------------------------------------------
\begin{frame}{Independencia: un problema muy común}

Supongamos que medimos:

\[
10\text{ hojas por planta}
\]

de:

\[
5\text{ plantas}
\]

Tenemos 50 mediciones, pero no necesariamente:

\[
n=50
\]

observaciones independientes.

Las hojas de una misma planta pueden estar correlacionadas.

\begin{alertblock}{Pseudorreplicación}
Tratar mediciones no independientes como réplicas independientes
puede producir inferencias engañosas.
\end{alertblock}

\end{frame}

% ------------------------------------------------------------
\begin{frame}{¿Qué hacer antes de una prueba?}

Antes de calcular un valor-$p$:

\begin{enumerate}
    \item Definir la pregunta científica.
    \item Identificar la unidad experimental.
    \item Revisar el diseño.
    \item Explorar los datos.
    \item Identificar variables y grupos.
    \item Considerar supuestos.
    \item Elegir el modelo/prueba.
\end{enumerate}

\end{frame}

% ------------------------------------------------------------
\begin{frame}{Flujo de trabajo}

\begin{center}

\begin{tikzpicture}[node distance=1.3cm]

\node[draw, rounded corners, fill=green!10] (q) {Pregunta biológica};
\node[draw, rounded corners, fill=blue!10, right=of q] (d) {Diseño};
\node[draw, rounded corners, fill=yellow!15, right=of d] (e) {Exploración};
\node[draw, rounded corners, fill=orange!15, right=of e] (t) {Prueba/modelo};
\node[draw, rounded corners, fill=red!10, right=of t] (i) {Interpretación};

\draw[->, thick] (q) -- (d);
\draw[->, thick] (d) -- (e);
\draw[->, thick] (e) -- (t);
\draw[->, thick] (t) -- (i);

\end{tikzpicture}

\end{center}

\vspace{0.5cm}

\centering
\textbf{La prueba estadística no sustituye el razonamiento científico.}

\end{frame}

% ------------------------------------------------------------
\begin{frame}[fragile]{Ejemplo en R}

Supongamos un conjunto de datos:

\begin{lstlisting}[language=R]
control <- c(12.1, 11.4, 13.2, 12.8, 11.9)
tratado <- c(15.0, 14.2, 16.1, 15.4, 14.8)

t.test(tratado, control)
\end{lstlisting}

La función produce, entre otros:

\begin{itemize}
    \item diferencia estimada,
    \item estadístico $t$,
    \item grados de libertad,
    \item valor-$p$,
    \item intervalo de confianza.
\end{itemize}

\end{frame}

% ------------------------------------------------------------
\begin{frame}[fragile]{Ejemplo en Python}

Con \texttt{scipy}:

\begin{lstlisting}[language=Python]
from scipy.stats import ttest_ind

control = [12.1, 11.4, 13.2, 12.8, 11.9]
tratado = [15.0, 14.2, 16.1, 15.4, 14.8]

resultado = ttest_ind(tratado, control)

print(resultado)
\end{lstlisting}

\end{frame}

% ------------------------------------------------------------
\begin{frame}{Interpretación de un resultado}

Supongamos que obtenemos:

\[
t=3.12
\]

\[
p=0.014
\]

\[
IC_{95\%}=[0.55,\;3.65]\text{ g}
\]

Una forma adecuada de comunicarlo:

\begin{block}
La biomasa media fue mayor en el grupo tratado que en el control,
con una diferencia estimada de 2.1 g (IC 95\%: 0.55--3.65 g;
$p=0.014$).
\end{block}

\end{frame}

% ------------------------------------------------------------
\begin{frame}{¿Significativo = importante?}

No necesariamente.

Consideremos:

\[
p<0.05
\]

Esto proporciona evidencia estadística contra $H_0$ según el análisis.

Pero todavía debemos preguntar:

\begin{itemize}
    \item ¿Qué tan grande es el efecto?
    \item ¿Es biológicamente relevante?
    \item ¿Cuál es su incertidumbre?
    \item ¿El diseño permite responder la pregunta?
    \item ¿Existen posibles sesgos?
\end{itemize}

\end{frame}

% ------------------------------------------------------------
\begin{frame}{Problema de comparaciones múltiples}

Supongamos que hacemos 20 pruebas independientes con:

\[
\alpha=0.05
\]

Incluso si todas las hipótesis nulas fueran ciertas,
la probabilidad de obtener al menos un resultado significativo puede ser
alta.

\[
P(\text{al menos un falso positivo})
=
1-(1-0.05)^{20}
\]

\[
\approx 0.64
\]

\begin{alertblock}{Idea}
Muchas pruebas requieren considerar el problema de multiplicidad.
\end{alertblock}

\end{frame}

% ------------------------------------------------------------
\begin{frame}{Correcciones por comparaciones múltiples}

Algunas estrategias:

\begin{itemize}
    \item Bonferroni
    \item Holm
    \item Benjamini--Hochberg
    \item Modelos que incorporan simultáneamente varios factores
\end{itemize}

La elección depende de:

\begin{itemize}
    \item objetivo inferencial,
    \item número de comparaciones,
    \item estructura del estudio,
    \item control deseado del error.
\end{itemize}

\end{frame}

% ------------------------------------------------------------
\begin{frame}{Ejercicio de discusión}

Un investigador compara la abundancia de aves entre dos bosques.

Obtiene:

\[
p=0.08
\]

Pregunta:

\begin{enumerate}
    \item ¿Podemos decir que ``no existe diferencia''?
    \item ¿Qué información adicional necesitamos?
    \item ¿Qué papel juega el tamaño de muestra?
    \item ¿Qué papel juega el tamaño del efecto?
    \item ¿Qué nos diría un intervalo de confianza?
\end{enumerate}

\vspace{0.3cm}

\pause

\begin{block}{Discusión}
No rechazar $H_0$ no equivale a demostrar que las poblaciones son idénticas.
\end{block}

\end{frame}

% ------------------------------------------------------------
\begin{frame}{Ejercicio: elegir la prueba}

Seleccione una estrategia para cada situación:

\begin{enumerate}
    \item Comparar crecimiento entre tratamiento y control.
    \item Comparar crecimiento antes y después en los mismos individuos.
    \item Comparar biomasa entre cuatro tratamientos.
    \item Evaluar asociación entre temperatura y tasa metabólica.
    \item Evaluar asociación entre dos variables categóricas.
\end{enumerate}

\vspace{0.3cm}

\pause

\textbf{Discusión:}

\begin{itemize}
    \item $t$ independiente
    \item $t$ pareada
    \item ANOVA
    \item correlación/regresión
    \item $\chi^2$ o Fisher
\end{itemize}

\end{frame}

% ------------------------------------------------------------
\begin{frame}{Ejercicio integrador}

Tenemos datos de crecimiento de plantas.

\begin{center}
\begin{tabular}{ccc}
\toprule
Planta & Control & Fertilizante\\
\midrule
1 & 10.2 & 12.1\\
2 & 11.4 & 12.8\\
3 & 9.8 & 11.9\\
4 & 10.7 & 13.0\\
5 & 10.9 & 12.5\\
6 & 11.2 & 13.1\\
\bottomrule
\end{tabular}
\end{center}

Preguntas:

\begin{enumerate}
    \item ¿Cuál es la unidad experimental?
    \item ¿Cuál es $H_0$?
    \item ¿Cuál es $H_A$?
    \item ¿Qué prueba usarían?
    \item ¿Qué supuestos revisarían?
    \item ¿Qué resultado reportarían?
\end{enumerate}

\end{frame}

% ------------------------------------------------------------
\begin{frame}{Ejercicio integrador: razonamiento}

Una posible formulación:

\[
H_0:\mu_F-\mu_C=0
\]

\[
H_A:\mu_F-\mu_C\neq0
\]

Si los grupos son independientes y los supuestos son razonables:

\[
\boxed{\text{prueba }t\text{ de dos muestras}}
\]

Pero la prueba final depende del diseño real.

\end{frame}

% ------------------------------------------------------------
\begin{frame}{Errores frecuentes}

\begin{itemize}
    \item Elegir la prueba mirando únicamente el valor-$p$.
    \item Interpretar $p$ como la probabilidad de que $H_0$ sea cierta.
    \item Confundir ``no significativo'' con ``sin efecto''.
    \item Ignorar la independencia.
    \item Hacer muchas pruebas sin considerar multiplicidad.
    \item Reportar únicamente el valor-$p$.
    \item Ignorar el tamaño del efecto.
    \item Convertir automáticamente resultados estadísticos en conclusiones
    biológicas.
\end{itemize}

\end{frame}

% ------------------------------------------------------------
\begin{frame}{¿Cómo reportar una prueba de hipótesis?}

Una buena comunicación incluye:

\begin{enumerate}
    \item Qué se comparó.
    \item Qué efecto se estimó.
    \item Tamaño del efecto.
    \item Intervalo de confianza.
    \item Prueba/modelo utilizado.
    \item Valor-$p$.
    \item Tamaño de muestra.
    \item Información relevante sobre el diseño.
\end{enumerate}

\vspace{0.3cm}

\begin{block}{Ejemplo}
El tratamiento aumentó la biomasa en 2.1 g
(IC 95\%: 0.55--3.65; $t=3.12$, $p=0.014$).
\end{block}

\end{frame}

% ------------------------------------------------------------
\begin{frame}{Mapa mental}

\begin{center}
\Large

\[
\boxed{\text{Pregunta biológica}}
\]

$\downarrow$

\[
\boxed{\text{Diseño y datos}}
\]

$\downarrow$

\[
\boxed{H_0\quad vs.\quad H_A}
\]

$\downarrow$

\[
\boxed{\text{Estadístico + incertidumbre}}
\]

$\downarrow$

\[
\boxed{\text{Valor-}p}
\]

$\downarrow$

\[
\boxed{\text{Tamaño del efecto}}
\]

$\downarrow$

\[
\boxed{\text{Interpretación biológica}}
\]

\end{center}

\end{frame}

% ------------------------------------------------------------
\begin{frame}{Resumen}

\begin{itemize}
    \item Las pruebas de hipótesis permiten cuantificar evidencia
    frente a una hipótesis nula bajo un modelo estadístico.
    \item El valor-$p$ no es la probabilidad de que una hipótesis sea
    verdadera.
    \item Debemos distinguir errores tipo I y tipo II.
    \item La potencia depende, entre otras cosas, del tamaño de muestra,
    efecto y variabilidad.
    \item La significancia estadística no equivale a importancia biológica.
    \item El diseño experimental y la independencia son fundamentales.
    \item Debemos reportar efectos e incertidumbre, no únicamente
    valores-$p$.
\end{itemize}

\end{frame}

% ------------------------------------------------------------
\begin{frame}{Preguntas para cerrar la clase}

\begin{enumerate}
    \item ¿Qué pregunta responde realmente un valor-$p$?
    \item ¿Cuál es la diferencia entre ``no rechazar $H_0$''
    y ``aceptar $H_0$''?
    \item ¿Qué es un error tipo I?
    \item ¿Qué factores afectan la potencia?
    \item ¿Por qué el tamaño del efecto es importante?
    \item ¿Por qué la independencia importa tanto en biología?
\end{enumerate}

\vspace{0.5cm}

\centering
\Large
\textbf{La estadística debe ayudar a responder preguntas biológicas,
no reemplazarlas.}

\end{frame}

% ------------------------------------------------------------
\begin{frame}{Actividad posterior}

\begin{block}{Tarea}
Tome un conjunto de datos biológicos y construya un análisis que incluya:
\end{block}

\begin{enumerate}
    \item Pregunta científica.
    \item Variable respuesta.
    \item Variable explicativa.
    \item Unidad experimental.
    \item $H_0$ y $H_A$.
    \item Visualización exploratoria.
    \item Prueba o modelo seleccionado.
    \item Tamaño del efecto.
    \item Intervalo de confianza.
    \item Valor-$p$.
    \item Interpretación biológica.
\end{enumerate}

\end{frame}

% ------------------------------------------------------------
\begin{frame}
\centering
\Huge
\textcolor{bioGreen}{Gracias}

\vspace{0.5cm}

\Large
¿Preguntas?
\end{frame}

% ============================================================
\end{document}
% ============================================================
